% Caractéristique intensité-tension et puissance d'un petit panneau photovoltaïque (L'essentiel, chapitre 25)
\begin{tikzpicture}
\begin{axis}[otMesures,otAxesBout,otLecture,width=100mm,height=52mm,xmin=0,xmax=6.6,ymin=0,ymax=3.8,
  xtick={0,1,...,6},ytick={0,1,...,3},
  xlabel={$U$ (V)},ylabel={$I$ (A) et $P$ (W)}]
\addplot[otNavy,line width=0.9pt,mark=*,mark size=1.3pt] coordinates {(0,0.85) (1,0.85) (2,0.84) (3,0.83) (4,0.80) (4.5,0.76) (5,0.66) (5.5,0.42) (6,0)};
\addplot[otRed!80!black,line width=0.9pt,mark=square*,mark size=1.2pt] coordinates {(0,0) (1,0.85) (2,1.68) (3,2.49) (4,3.20) (4.5,3.42) (5,3.30) (5.5,2.31) (6,0)};
\draw[otGreenDark,line width=0.5pt,dash pattern=on 1.5pt off 1pt] (axis cs:4.5,0) -- (axis cs:4.5,3.42);
\node[pcLegende,anchor=south,text=otGreenDark] at (axis cs:4.5,3.45) {$P_\mathrm{max}=\qty{3,42}{W}$};
\node[pcLegende,anchor=south west,text=otNavy] at (axis cs:0.2,0.9) {intensité $I$};
\node[pcLegende,anchor=south east,text=otRed!80!black] at (axis cs:2.8,2.5) {puissance $P=UI$};
\end{axis}
\end{tikzpicture}
\par\smallskip\legende{La puissance électrique fournie passe par un maximum ; le rendement se calcule à ce point de fonctionnement.}
